\chapter{Introduction} \label{chap:intro} \section{Clinical background and motivation} Sepsis-associated acute kidney injury (SA\mbox{-}AKI) has emerged as one of the most consequential syndromes in critical care because it couples the systemic instability of sepsis with acute loss of renal function and thereby compounds risk across organ systems. For clinical and research purposes SA\mbox{-}AKI is commonly labeled when acute kidney injury develops within seven days of sepsis onset so that kidney dysfunction can be attributed to the septic process rather than to unrelated peri-ICU events \citep{Peerapornratana_2019}. The constituent conditions are defined by international consensus. Sepsis is defined by Sepsis\mbox{-}3 as life-threatening organ dysfunction due to a dysregulated host response to infection and is operationalized by an acute rise in SOFA score of two points or more \citep{Singer_2016}. Acute kidney injury is defined and staged by KDIGO using serum creatinine and urine output thresholds, including an absolute creatinine increase of at least $0.3\,\mathrm{mg/dL}$ within forty-eight hours, a relative increase of at least $1.5\times$ baseline within seven days, or oliguria at or below $0.5\,\mathrm{mL/kg/h}$ for six hours \citep{KDIGO_2012}. Within intensive care units SA\mbox{-}AKI has been reported to account for a large share of AKI cases and to be associated with sharply higher in-hospital mortality, longer length of stay, and increased risk of subsequent chronic kidney disease in survivors \citep{Poston_Koyner_2019,Peerapornratana_2019}. These epidemiological features have kept early risk stratification at the forefront of research because timely identification of clinical deterioration is central to appropriate escalation, nephroprotective strategies, and resource planning. \section{What the literature has tried and where it falls short} Risk assessment in the ICU has traditionally relied on general severity scores such as SOFA, SAPS~II, and APACHE~II. These scores have offered useful global snapshots of illness severity but they tend to show only fair discrimination when the target is SA\mbox{-}AKI mortality, and they do not represent the nonlinear and interacting physiology that links infection biology, hemodynamics, renal injury, and treatment exposure in this population \citep{Luo_2022,Qu_2024}. In recent years the availability of large electronic health record repositories, notably MIMIC\mbox{-}IV, has encouraged the development of machine learning models that can take advantage of richer feature spaces and flexible function classes \citep{Johnson_2023_MIMIC}. Tree-based ensembles such as extreme gradient boosting frequently outperform baseline scores in internal testing and AUROC values near or above 0.80 are often reported \citep{Chen_2025,Qu_2024,Wang_2024}. Nevertheless, several methodological patterns recur across studies and create uncertainty about the true clinical utility of published models. Temporal anchoring of predictors is sometimes vague and whole-stay aggregation has been used, which allows future information to flow into features that are later presented as early predictors. Proxy variables for time to event, including length of stay, have been included as inputs to mortality models, which introduces a direct leakage channel from outcome process to covariates and inflates apparent discrimination \citep{Chen_2025}. Data-dependent steps such as feature selection and imputation have sometimes been performed across the full dataset rather than being learned on development folds and then applied to held-out data, which again invites optimistic estimates \citep{Kang_Yoon_2023}. Class imbalance has often been handled through synthetic oversampling, including SMOTE and ADASYN, without careful assessment of whether the generated examples concentrate in overlapped or noisy regions or near outliers where decision boundaries are unstable \citep{Chawla_2002,He_2008,Blagus_2013,Saez_2015,Krawczyk_2016}. Calibration and principled threshold selection have been inconsistently reported, which limits bedside interpretability even when AUROC values appear strong. Finally, most model development has relied on data from a single tertiary-care center, and performance losses under transport to other hospitals and populations have been documented, which tempers claims of broad generalizability \citep{Johnson_2023_MIMIC,Rockenschaub_2024}. \section{Problem focus and research objectives} The central challenge addressed in this thesis is the chasm between the promising internal accuracy reported in machine learning literature and the practical bedside utility of predictive models for SA-AKI mortality. This gap is primarily shaped by two issues: first, a failure to align predictive pipelines with the real-world clinical decision timeline, and second, a lack of transparency and calibration that hinders a model's adoption. The objective of this work is to develop and validate a mortality risk stratification framework for SA\mbox{-}AKI that is methodologically rigorous, clinically relevant, and transparent. The work is conducted using the MIMIC\mbox{-}IV dataset \citep{Johnson_2023_MIMIC}, with a strict focus on an early clinical horizon to ensure the model's predictions are actionable. The specific aims are: \begin{enumerate} \item To perform a systematic exploratory data analysis to identify key clinical phenotypes and risk factors within the training cohort. \item To develop and validate a high-performing, calibrated LightGBM classification model for in-hospital mortality using only features available within the first 24 hours of ICU admission. \item To ensure model transparency by providing both global and local feature-level explanations using SHapley Additive exPlanations (SHAP) \citep{Lundberg_Lee_2017}. \item To complement the binary prediction with a dynamic risk assessment by developing survival models using Cox Proportional Hazards and DeepSurv to estimate the \emph{time to in-hospital death}. \end{enumerate} \section{Proposed approach and evaluation plan} Our approach begins with the disciplined construction of an SA\mbox{-}AKI cohort according to Sepsis\mbox{-}3 and KDIGO criteria, with all predictors strictly limited to the first 24 hours (T0--T24) of an ICU stay \citep{Singer_2016,KDIGO_2012}. This work is distinguished by a robust, leakage-aware preprocessing pipeline where all steps are fitted exclusively on training data. This includes a formal test for the missing data mechanism, followed by the creation and statistical selection of informative missingness indicators and median imputation. A comprehensive benchmark of machine learning models is conducted, leading to the selection of a final LightGBM model based on its superior performance. This model is meticulously tuned using the Optuna framework, explicitly calibrated using isotonic regression, and an optimal decision threshold is determined via Youden's J statistic on a validation set. Explanations for the final model's predictions are generated using SHAP to ensure transparency. Finally, to provide a dynamic view of risk over time, we develop and evaluate two survival models—a traditional Cox Proportional Hazards model and a neural network-based DeepSurv model—using the same early-window covariates. The evaluation plan prioritizes methodological transparency and clinical relevance, reporting a suite of metrics including AUROC, F1-score, Harrell's C-index, and calibration plots on a strictly held-out test set, addressing the common pitfalls identified in the literature \citep{Kang_Yoon_2023,Chen_2025}. \section{Chapter roadmap} The remainder of the thesis is organized to follow this arc from clinical framing to deployment-oriented evaluation. Chapter~\ref{chap:litreview} reviews the background in SA\mbox{-}AKI epidemiology and definitions, surveys conventional scores and machine learning approaches, and synthesizes methodological lessons from recent studies. Chapter~\ref{chap:problem} states the formal problem and gives the mathematical formulation for the binary prediction and survival analysis tasks. Chapter~\ref{chap:methods} provides a detailed description of the cohort construction, the leakage-aware preprocessing pipeline, the model development and calibration process, and the evaluation protocol. Chapter~\ref{chap:eda} presents the findings from the extensive exploratory data analysis of the training cohort, identifying key clinical phenotypes and variable distributions that inform the modeling process. Chapter~\ref{chap:results} presents the final performance results for the classification and survival models on the held-out test set, including discrimination, calibration, and SHAP-based interpretations. Finally, Chapter~\ref{chap:conclusion} concludes with a summary of findings, comments on clinical implications, and outlines future research directions.